\section{Design study on development datasets}
\label{sec:dev}

Before touching CC18 we evaluated eleven candidate changes to NEXUS~v1 on 16
development datasets (Australian, jasmine, sylvine, Titanic, hill-valley,
monks-problems-2, wine-quality-red, ionosphere, sonar, yeast, page-blocks,
glass, banana, pc2, waveform-5000, heart-statlog; 155--8{,}378 rows, 2--10
classes) with 5-fold cross-validation. Member predictions were cached so that all
candidates are compared on identical base models. \Cref{tab:dev} lists every
candidate, including those that hurt.

\begin{table}[t]
\centering
\caption{Every design candidate evaluated on the 16 development datasets,
relative to NEXUS~v1. $\Delta R^2$: mean reduction in log loss as a fraction of
label entropy (percentage points). Rel.\ LL / Rel.\ err.: mean relative
reduction of log loss / classification error. Wins: datasets where the candidate
has lower log loss than v1. The bold row was frozen as NEXUS~v2.}
\label{tab:dev}
\small
\begin{tabular}{lrrrrr}
\toprule
Candidate & Rank & $\Delta R^2$ (pp) & Rel.\ LL & Wins/16 & Rel.\ err.\\
\midrule
+extended portfolio +L2 stack (3) & 4.88 & +0.64 & +7.6\% & 10 & +6.5\%\\
+L2 stack \{lgbm, logreg, catboost\} & 5.34 & +0.64 & +7.5\% & 9 & +7.6\%\\
\textbf{+L2 stack \{lgbm, logreg\} (v2)} & \textbf{6.03} & \textbf{+0.62} & \textbf{+7.0\%} & \textbf{11} & \textbf{+7.2\%}\\
+L2 stack \{lgbm, catboost\} & 5.09 & +0.48 & +7.2\% & 11 & +6.9\%\\
+L2 stack \{lgbm\} & 6.09 & +0.38 & +6.3\% & 12 & +7.1\%\\
repeated cross-fitting R=3 & 7.12 & +0.05 & -1.5\% & 8 & -1.1\%\\
repeated cross-fitting R=2 & 6.44 & +0.04 & -0.9\% & 11 & -0.9\%\\
+13-config extended portfolio & 6.03 & +0.03 & +0.1\% & 9 & -0.6\%\\
v1 (greedy selection + temperature) & 7.34 & +0.00 & +0.0\% & 0 & +0.0\%\\
convex weight optimisation & 7.44 & -0.04 & -0.0\% & 8 & -0.3\%\\
early-stopped GBDTs (swap) & 7.75 & -0.22 & -5.3\% & 7 & -1.7\%\\
bagged ensemble selection & 8.44 & -0.49 & -10.5\% & 5 & -0.9\%\\
\bottomrule
\end{tabular}

\end{table}

\paragraph{What helped.} Only a second stacking layer gave a consistent gain
(\Cref{tab:dev}). Using LightGBM and logistic regression as layer-2 models
captured 97\% of the gain of the best three-model stack at about one ninth of
its additional cost (CatBoost as a layer-2 model alone cost more than all eight
layer-1 models together), so this configuration became NEXUS~v2.

\paragraph{What did not.} A 13-configuration extended portfolio (early-stopped
GBDTs of several sizes, randomised LightGBM, entropy-split forests, alternative
$k$-NN and MLP settings) added almost nothing on top of the eight-model library,
even though its early-stopped CatBoost was the single best model. Replacing the
three GBDTs with early-stopped versions, bagged ensemble selection
\citep{caruana2006getting} and exact convex optimisation of the weights did not
improve on plain greedy selection.

\paragraph{A systematic bias in out-of-fold selection.} Selection uses
single-model OOF predictions, but test-time predictions average $K$ fold models.
\Cref{tab:bagging} shows that the improvement from this averaging differs
substantially between algorithms, from about 4\% for logistic regression to
15\% for $k$-NN, and is consistent across datasets. OOF estimates therefore
understate the test-time quality of high-variance learners more than that of
low-variance ones. The principled remedy, repeated cross-fitting so that each
OOF prediction is itself an average, did not improve the final ensemble on
average (\Cref{tab:dev}): with eight diverse members, greedy selection appears
robust to this bias. We report it as a measured effect and a negative result.

\begin{table}[t]
\centering
\caption{Bagging gain per algorithm on the development datasets: mean OOF log
loss of single fold models versus mean test log loss of the 5-model bag, and
the mean per-dataset relative improvement.}
\label{tab:bagging}
\small
\begin{tabular}{lrrrr}
\toprule
Algorithm & OOF LL & Bagged test LL & Gain (mean) & Datasets with gain\\
\midrule
Logistic regression & 0.498 & 0.484 & 4.3\% & 12/16\\
MLP & 0.517 & 0.492 & 4.7\% & 16/16\\
XGBoost & 0.444 & 0.416 & 5.5\% & 15/16\\
Random forest & 0.430 & 0.394 & 7.4\% & 15/16\\
CatBoost & 0.401 & 0.365 & 7.8\% & 14/16\\
Extra trees & 0.424 & 0.390 & 8.5\% & 15/16\\
LightGBM & 0.543 & 0.465 & 11.9\% & 14/16\\
k-NN & 0.828 & 0.681 & 14.8\% & 16/16\\
\bottomrule
\end{tabular}

\end{table}
